\chapter{Conclusiones}

Las conclusiones sintetizan, de forma clara y ordenada, la respuesta al planteamiento planteado en el Capítulo 1 y a los objetivos fijados en el desarrollo. Se prioriza la vinculación con los resultados y la contribución específica, evitando repetir de forma extensa el estado del arte o la descripción de métodos.

\section{Respuesta al objetivo general y a los específicos}
Sintetice en qué medida se ha logrado el encaje entre modelado generativo, teoría de la información y criterio de riesgo basado en la divergencia de Kullback-Leibler, en relación con su aplicación a la optimización de carteras. Si fijó objetivos parcialmente cumplidos, señálese con precisión.

\section{Contribuciones y posición frente a la literatura}
Destaque las aportaciones propias (técnicas, de diseño experimental, de criterio de evaluación) y cómo se sitúan respecto a los enfoques revisados, sin reescribir el capítulo 2, en un máximo de tono conclusivo y prospectivo mínima.

\section{Síntesis de hallazgos}
Puede ofrecerse un breve cierre que recoja, en lenguaje accesible, los principales hallazgos cuantitativos o metodológicos, siempre que se eviten afirmaciones no respaldadas por el cuerpo de la memoria.
